% ============================================================
% PROYECTO INTEGRADOR: ODOO PARA LA GESTIÓN DEL PAÑOL
% ============================================================
\clearpage
\section*{PROYECTO INTEGRADOR}
\subsection*{Implementación de un servidor Odoo para la gestión del pañol}

\begin{recuadro}
\textbf{Consigna general:} implementar en la escuela un servidor con \textbf{Odoo Community} que
gestione el \textbf{pañol} de la especialidad: herramientas, instrumental, stock de componentes
electrónicos, almacenes y ubicaciones, \textbf{pedidos de materiales}, \textbf{órdenes de
trabajo} y \textbf{préstamos} de herramientas, aplicando lo visto en los módulos de 5S/Lean y
ERP.

\textbf{Modalidad:} trabajo en equipos con roles rotativos. \textbf{Entrega final:} servidor
funcionando con datos reales, manual de usuario y presentación oral.
\end{recuadro}

% ------------------------------------------------------------
\section*{1. ORGANIZACIÓN DEL CURSO}

\begin{center}
\small
\begin{tabular}{|p{0.27\linewidth}|p{0.63\linewidth}|}
\hline
\textbf{Equipo} & \textbf{Responsabilidad} \\ \hline
Infraestructura & Instalación del servidor, base de datos, copias de seguridad, red y usuarios. \\ \hline
Relevamiento y 5S & Inventario físico, clasificación, codificación y etiquetado del pañol. \\ \hline
Datos maestros & Carga de productos, categorías, ubicaciones, proveedores y reglas de
reposición. \\ \hline
Procesos & Configuración de pedidos de materiales, préstamos y órdenes de trabajo. \\ \hline
Documentación y capacitación & Manual de usuario, procedimientos y capacitación a docentes y
pañoleros. \\ \hline
\end{tabular}
\end{center}

% ------------------------------------------------------------
\section*{2. ETAPA 1: RELEVAMIENTO DEL PAÑOL (5S)}

Antes de cargar un solo dato en el sistema, el pañol tiene que estar ordenado: un ERP con datos
desordenados solo produce informes desordenados.

\begin{enumerate}
  \item \textbf{Seiri:} separar lo que está roto, obsoleto o no pertenece al pañol.
  \item \textbf{Seiton:} definir el esquema de ubicaciones y etiquetar estanterías y cajones.
  \item \textbf{Relevamiento:} planilla con código, descripción, cantidad, ubicación, estado y
        si se controla por número de serie.
  \item \textbf{Codificación:} definir una regla para los códigos internos. Ejemplo:
\end{enumerate}

\begin{center}
\small
\begin{tabular}{|l|l|l|}
\hline
\textbf{Prefijo} & \textbf{Categoría} & \textbf{Ejemplo} \\ \hline
HER- & Herramientas de mano & HER-0012 Pinza de punta \\ \hline
INS- & Instrumental (con n° de serie) & INS-0003 Osciloscopio 100 MHz \\ \hline
CMP- & Componentes electrónicos & CMP-RES-10K Resistor 10 k$\Omega$ 1/4 W \\ \hline
CON- & Consumibles & CON-0005 Estaño 0,8 mm \\ \hline
CAB- & Cables y conectores & CAB-0021 Conector BNC macho \\ \hline
\end{tabular}
\end{center}

% ------------------------------------------------------------
\section*{3. ETAPA 2: INSTALACIÓN DEL SERVIDOR}

\subsubsection*{3.1. Requisitos}
\begin{itemize}
  \item Una PC o mini-PC dedicada (4 GB de RAM como mínimo, 8 GB recomendado, disco SSD) con
        Linux (Ubuntu Server o Debian), conectada a la red de la escuela.
  \item Alternativa para pruebas: cualquier PC con \textbf{Docker}.
  \item IP fija en la red local para acceder desde el navegador:
        \texttt{http://IP-DEL-SERVIDOR:8069}.
\end{itemize}

\subsubsection*{3.2. Instalación con Docker Compose}
Se crea una carpeta \texttt{odoo-panol} con el archivo \texttt{docker-compose.yml}:

\begin{verbatim}
services:
  db:
    image: postgres:16
    environment:
      - POSTGRES_USER=odoo
      - POSTGRES_PASSWORD=CAMBIAR_CLAVE
      - POSTGRES_DB=postgres
    volumes:
      - db-data:/var/lib/postgresql/data
    restart: always

  odoo:
    image: odoo:19
    depends_on:
      - db
    ports:
      - "8069:8069"
    environment:
      - HOST=db
      - USER=odoo
      - PASSWORD=CAMBIAR_CLAVE
    volumes:
      - odoo-data:/var/lib/odoo
      - ./addons:/mnt/extra-addons
    restart: always

volumes:
  db-data:
  odoo-data:
\end{verbatim}

Y se levanta con:
\begin{verbatim}
sudo docker compose up -d
\end{verbatim}

Al ingresar por primera vez a \texttt{http://IP-DEL-SERVIDOR:8069} se crea la base de datos:
nombre (\texttt{panol\_eest7}), correo y contraseña del administrador, idioma
\textbf{Español (AR)} y país \textbf{Argentina}. \textbf{No} marcar ``datos de demostración''
en la base real (sí en una base de práctica).

\subsubsection*{3.3. Aplicaciones a instalar}
\begin{itemize}
  \item \textbf{Inventario:} productos, almacenes, ubicaciones, movimientos, números de serie.
  \item \textbf{Compras:} solicitudes de presupuesto y órdenes de compra a proveedores.
  \item \textbf{Mantenimiento:} equipos (instrumental) y solicitudes de mantenimiento o
        calibración.
  \item \textbf{Reparaciones:} órdenes de trabajo que consumen componentes del stock.
  \item \textbf{Empleados} y \textbf{Contactos:} docentes, pañoleros, estudiantes y proveedores.
\end{itemize}

\subsubsection*{3.4. Seguridad y copias de respaldo}
\begin{itemize}
  \item Cambiar todas las contraseñas de ejemplo y la \textbf{contraseña maestra} del
        administrador de bases de datos.
  \item Crear usuarios con \textbf{permisos mínimos}: administrador (docente), pañolero
        (inventario completo), estudiante (solo crear pedidos).
  \item \textbf{Copia de seguridad semanal} desde el gestor de bases de datos
        (\texttt{/web/database/manager}) guardada fuera del servidor.
  \item No exponer el servidor a Internet sin el asesoramiento del área de informática.
\end{itemize}

% ------------------------------------------------------------
\section*{4. ETAPA 3: CONFIGURACIÓN DE ALMACENES Y PRODUCTOS}

\subsubsection*{4.1. Ajustes de Inventario}
En \textit{Inventario $\rightarrow$ Configuración $\rightarrow$ Ajustes} activar:
\textbf{Ubicaciones de almacenamiento}, \textbf{Rutas de varios pasos} (opcional),
\textbf{Lotes y números de serie} y \textbf{Unidades de medida}.

\subsubsection*{4.2. Almacenes y ubicaciones}
\begin{verbatim}
PAN (Almacén Pañol)
 |-- PAN/Stock/Est-A/Caj-01 ... Caj-20   (componentes)
 |-- PAN/Stock/Tablero-Herramientas      (herramientas de mano)
 |-- PAN/Stock/Armario-Instrumental      (instrumental)
 |-- PAN/Stock/Consumibles
TAL (Almacén Taller de Aviónica)
 |-- TAL/Stock/Banco-01 ... Banco-08     (puestos de trabajo)
Ubicaciones virtuales:
 |-- Préstamos/Estudiantes
 |-- Bajas/Rotos
\end{verbatim}

\subsubsection*{4.3. Productos}
\begin{itemize}
  \item \textbf{Componentes y consumibles:} tipo \textit{bienes}, con seguimiento de inventario,
        código interno, categoría, unidad de medida y proveedor.
  \item \textbf{Instrumental y herramientas valiosas:} seguimiento \textbf{por número de
        serie}, para saber en todo momento dónde está cada equipo y quién lo tiene.
  \item Cargar los datos con la \textbf{importación desde planilla} (archivo CSV o XLSX con la
        plantilla que ofrece Odoo), a partir del relevamiento de la etapa 1.
  \item Cargar el stock inicial mediante \textbf{Inventario físico} (ajuste de inventario).
  \item Imprimir \textbf{etiquetas con código de barras} para cajones y equipos.
\end{itemize}

\subsubsection*{4.4. Reglas de reabastecimiento (Kanban digital)}
Para cada componente de alto consumo se define una regla mínimo/máximo. Ejemplo: Resistor 10
k$\Omega$, mínimo 50 unidades, máximo 200. Cuando el stock baja de 50, Odoo propone una
\textbf{solicitud de presupuesto} al proveedor por la cantidad necesaria para llegar a 200.

% ------------------------------------------------------------
\section*{5. ETAPA 4: PROCESOS DEL PAÑOL}

\subsubsection*{5.1. Pedido de materiales}
Se crea un \textbf{tipo de operación} ``Pedido de materiales'' (transferencia interna) desde
\texttt{PAN/Stock} hacia \texttt{TAL/Stock}.
\begin{enumerate}
  \item El estudiante o docente carga el pedido: componentes, cantidades y banco de destino.
  \item El pañolero lo revisa, prepara los materiales y \textbf{valida} la transferencia.
  \item El stock del pañol se descuenta automáticamente y queda registrado quién lo pidió, qué
        se entregó y cuándo.
\end{enumerate}

\subsubsection*{5.2. Préstamo y devolución de herramientas}
\begin{enumerate}
  \item Préstamo: transferencia interna desde \texttt{PAN/Stock} hacia
        \texttt{Préstamos/Estudiantes}, indicando el estudiante como contacto y el número de
        serie del equipo.
  \item Devolución: transferencia inversa. Si vuelve dañado, se envía a \texttt{Bajas/Rotos}
        o se genera una solicitud de mantenimiento.
  \item Al final de la clase, el informe de stock en \texttt{Préstamos/Estudiantes} muestra
        qué falta devolver (control de herramientas y prevención de FOD).
\end{enumerate}

\subsubsection*{5.3. Órdenes de trabajo}
Con la aplicación \textbf{Reparaciones}, cada trabajo del taller se registra como una
\textbf{orden de reparación}:
\begin{itemize}
  \item Equipo a reparar (por ejemplo, un transceptor VHF de práctica) y descripción de la
        falla.
  \item Técnico responsable y estado (borrador, confirmada, en curso, terminada).
  \item \textbf{Componentes utilizados}, que se descuentan del stock al finalizar la orden.
  \item Notas del trabajo realizado y horas empleadas.
\end{itemize}
Así cada orden queda con su \textbf{costo de materiales} y su \textbf{trazabilidad}, igual que
en un taller aeronáutico habilitado.

\subsubsection*{5.4. Mantenimiento y calibración del instrumental}
Con la aplicación \textbf{Mantenimiento} se dan de alta los equipos (osciloscopios, fuentes,
multímetros, estaciones de soldado) con su número de serie, y se programan
\textbf{mantenimientos preventivos} periódicos (por ejemplo, verificación anual de
calibración). Las fallas se informan como \textbf{solicitudes de mantenimiento correctivo}.

\subsubsection*{5.5. Compras}
Desde las propuestas de reposición se generan las \textbf{solicitudes de presupuesto}, se
comparan proveedores, se confirma la \textbf{orden de compra} y, al llegar la mercadería, se
\textbf{valida la recepción}: el stock aumenta y la orden queda cerrada.

% ------------------------------------------------------------
\section*{6. ETAPA 5: INFORMES Y MEJORA CONTINUA}

\begin{itemize}
  \item \textbf{Valorización del inventario:} cuánto dinero hay inmovilizado en el pañol.
  \item \textbf{Componentes más consumidos} (análisis de Pareto) para ajustar los mínimos.
  \item \textbf{Herramientas sin devolver} y equipos fuera de servicio.
  \item \textbf{Órdenes de trabajo por estado} y tiempo promedio de resolución.
  \item \textbf{Propuesta Kaizen:} con los datos de los informes, cada equipo propone una mejora
        del proceso y la presenta con el ciclo PDCA.
\end{itemize}

% ------------------------------------------------------------
\section*{7. ENTREGABLES Y EVALUACIÓN}

\begin{center}
\small
\begin{tabular}{|p{0.62\linewidth}|p{0.25\linewidth}|}
\hline
\textbf{Entregable} & \textbf{Peso} \\ \hline
Relevamiento del pañol, codificación y fotos antes/después (5S) & 15\% \\ \hline
Servidor instalado, usuarios, permisos y copias de seguridad & 20\% \\ \hline
Datos maestros cargados (productos, ubicaciones, reglas de reposición) & 20\% \\ \hline
Procesos funcionando: pedido de materiales, préstamos, órdenes de trabajo, compras & 25\% \\ \hline
Manual de usuario y procedimientos (con capturas de pantalla) & 10\% \\ \hline
Presentación oral y propuesta de mejora (Kaizen) & 10\% \\ \hline
\end{tabular}
\end{center}

\begin{recuadro}
\textbf{Para el CV:} al terminar el proyecto, cada estudiante puede incluir en su currículum y
en LinkedIn la experiencia de \textit{``Implementación de ERP Odoo para la gestión de
inventario, compras y órdenes de trabajo de un pañol técnico''}, junto con las tareas
concretas que realizó en su equipo.
\end{recuadro}
